\section*{Chapter \ref{ch:in:revenue-and-profit-across-the-industry} --- Revenue and Profit across the Industry}

\begin{solution}{exo:in:revenue-and-profit-across-the-industry:1}
Revenues less transaction-based expenses are operating expenses plus operating income, $5\,002+4\,929=\$9\,931$ million;
the transaction-based expenses are $12\,640-9\,931=\$2\,709$ million.
\end{solution}

\begin{solution}{exo:in:revenue-and-profit-across-the-industry:2}
Nasdaq (\$0.55 million), Man Group and the Intercontinental Exchange (\$0.77 million each), Morgan Stanley (\$0.85
million), Flow Traders (\$0.86 million) and MarketAxess (\$0.94 million). The highest whole-firm figure, Quadrature's
\$9.29 million, is 16.9 times Nasdaq's.
\end{solution}

\begin{solution}{exo:in:revenue-and-profit-across-the-industry:3}
CME Group: $4\,229.5/3\,875=\$1.09$ million. Optiver: $1\,369/2\,100=\euro0.65$ million.
\end{solution}

\begin{solution}{exo:in:revenue-and-profit-across-the-industry:4}
Its headcount rose from 6\,377 to 8\,525 during 2023, ``primarily due to our acquisition of Adenza'': the acquired staff
count in full at the year end while the acquired revenue arrives over the following years.
\end{solution}

\begin{solution}{exo:in:revenue-and-profit-across-the-industry:5}
$(29.25/16.85)^{1.19}=1.92$: \omterm{def:in:revenue-and-profit-across-the-industry:perhead}{revenue per head} about 92\% higher. Virtu's actual ratio was 2.16 and Flow Traders' 4.51:
the fitted elasticity averages two firms whose sensitivity differs, and a year's mean VIX is one of several things that
set a market maker's revenue.
\end{solution}

\begin{solution}{exo:in:revenue-and-profit-across-the-industry:6}
It is a subsidiary whose revenue includes service income set by the group's transfer pricing and whose headcount fell
from 96 to 21 as work moved elsewhere in the group: its \omterm{def:in:revenue-and-profit-across-the-industry:perhead}{revenue per head} (\$1.00 million to \$7.18 million) describes the
entity's accounts, not the firm's people.
\end{solution}

\begin{solution}{exo:in:revenue-and-profit-across-the-industry:7}
The market makers' elasticity stays 1.19 and the exchanges' 0.13; their standard errors change only in the third
decimal. With firm effects and one slope per kind, each kind's slope is estimated from
its own firms' variation; the other kinds enter only through the pooled residual variance.
\end{solution}

\begin{solution}{exo:in:revenue-and-profit-across-the-industry:8}
The table's top is the top of those who file, not of the industry: the largest private market makers publish no
consolidated income, so \omterm{def:in:revenue-and-profit-across-the-industry:disclosure}{disclosure bias} removes them. The \$13 million is also a different line (a fund manager's fees)
from a market maker's net trading revenue, and it comes from 98 people in one exceptional year.
\end{solution}

\begin{solution}{pb:in:revenue-and-profit-across-the-industry:1}
\begin{enumerate}
\item Revenue (after pass-through costs) or a stated profit line, divided by headcount; with the line, the headcount's
basis, the currency and the price level.
\item A gross line counts money collected for others; the exchange group's total revenues of \$12\,640 million include
\$2\,709 million of transaction-based expenses, leaving \$9\,931 million.
\item Year end, average, full-time equivalents, rounded thousands, stated lower bounds; year-end counts overstate the
denominator for a growing firm, and a lower bound gives an upper bound per head.
\item ECB annual average rates to dollars, then the US consumer price index to 2025 dollars.
\item Fifteen whole firms (plus one UK entity shown apart), 87 firm-years, from SEC company facts, Book~16's tables,
annual reports, results releases and the UK register.
\item Lowest: Nasdaq, \$0.48 million in 2023. Highest: Quadrature, \$13.0 million in its year to January 2022.
\item Flow Traders 4.5, Virtu 2.5, CME Group 1.2.
\item Acquisitions: Nasdaq's headcount rose from 6\,377 to 8\,525 with Adenza; the Intercontinental Exchange's from
8\,911 to 13\,222.
\item Banks file markets revenue but no division headcount.
\item It rose from \$1.00 million to \$7.18 million as its \omterm{def:in:reading-a-trading-firms-accounts:staff}{average headcount} fell from 96 to 21; its revenue is set within
the group.
\item 15.39, 29.25, 19.66, 25.64, 16.85, 15.55, 18.93.
\item $\log(R/N)_{it}=a_i+b_{k(i)}\log\bar V_t+\varepsilon_{it}$; the firm effects remove each firm's level, so the
slopes come from movements within firms.
\item Market makers 1.19 (0.26); exchanges and venues 0.13 (0.16); asset managers 0.23 (0.27); banks 0.03 (0.32);
broker $-0.39$ (0.37); crypto exchange $-1.27$ (0.59).
\item Its revenue follows crypto prices, which fell hardest in 2022, a year of high equity volatility.
\item Virtu 2.16, Flow Traders 4.51.
\item Market makers 1.19 (0.26), exchanges and venues 0.13 (0.16), asset managers 0.23 (0.27), banks 0.03 (0.32),
broker $-0.39$ (0.37), crypto exchange $-1.27$ (0.59); the 2020/2023 ratio of \omterm{def:in:revenue-and-profit-across-the-industry:perhead}{revenue per head} ranges from 2.16 to 4.51
across the market makers.
\item The residuals share years and are not independent; seven years hold one large volatility event.
\item The bias from the publishing firms differing from those that do not; it leaves out the most profitable private
firms, so it understates the top.
\item Upward: firms that failed leave with their bad years.
\item Yes, the market maker has good years: its revenue per person roughly doubled or more from 2023 to 2020, and it moves
with the market far more than an exchange's. The exchange's is steadier and the market maker's variable pay will be too
(chapter~13).
\end{enumerate}
\end{solution}

\begin{solution}{iq:in:revenue-and-profit-across-the-industry:1}
More volume (more trades to earn a spread on) and wider spreads (more earned per trade, as inventory risk rises); also
more dislocations between related instruments.
\iqlookfor{volume and spread, with the inventory-risk reason for the spread.}
\end{solution}

\begin{solution}{iq:in:revenue-and-profit-across-the-industry:2}
Revenues less transaction-based expenses, because the rest is collected for others (regulatory fees) or paid back to
liquidity providers; a bank's net revenues are likewise after interest expense.
\iqlookfor{netting pass-through costs before comparing.}
\end{solution}

\begin{solution}{iq:in:revenue-and-profit-across-the-industry:3}
The firm effect removes each firm's permanent level, so the slope comes from within-firm movement. A year effect would
absorb every common year shock, including the VIX itself, which varies only by year: the slope would no longer be
identified, except as differences between kinds.
\iqlookfor{identification: a year-level regressor cannot coexist with year effects.}
\end{solution}

\begin{solution}{iq:in:revenue-and-profit-across-the-industry:4}
Cluster by year (the shock is common), but seven clusters are too few for the usual formula; use a wild bootstrap by
year, or report the range of slopes firm by firm.
\iqlookfor{few clusters, and a method that respects them.}
\end{solution}

\begin{solution}{iq:in:revenue-and-profit-across-the-industry:5}
Key each fact by (filer, concept, period start, period end, filing); keep every filing's value so restatements are new
rows; a view picks the latest filing per period. Store the unit and currency with the value.
\iqlookfor{bitemporal keys: period and filing.}
\end{solution}

\begin{solution}{iq:in:revenue-and-profit-across-the-industry:6}
From public pieces only: headcount from their own statements (as ranges), revenue ranges from the listed peers'
\omterm{def:in:revenue-and-profit-across-the-industry:perhead}{revenue per head} scaled by any public activity measure, stated as a range with its assumptions. I would refuse to
present it as a number, to rank firms by it, or to use press figures as data.
\iqlookfor{ranges, stated assumptions, and refusing false precision.}
\end{solution}
